% =============================================================
%  Module Description — Data Structures and Algorithms I (CS201-S27)
%  DRAFT for a planned module (Semester 2 of the Foundation programme).
%  The code CS201-S27 is a proposal (semester digit + position in the
%  semester); add the term suffix (e.g. -S27) once the term is known and
%  make it equal the course's `code` in src/content/courses/.
%  Compile with: pdflatex cs201-module-description.tex
%  Write "TODO(owner): ..." where information is still missing;
%  the website highlights it.
% =============================================================
\documentclass[11pt,a4paper]{article}

\usepackage[utf8]{inputenc}
\usepackage[T1]{fontenc}
\usepackage{lmodern}
\usepackage[margin=2.2cm]{geometry}
\usepackage{booktabs}
\usepackage{tabularx}
\usepackage{array}
\usepackage{enumitem}
\usepackage{xcolor}
\usepackage{titlesec}
\usepackage{fancyhdr}
\usepackage[hidelinks]{hyperref}

% ---------- Module details (edit these) ----------------------
% \ModuleCode must equal the course's `code` in src/content/courses/.
\newcommand{\ModuleTitle}{Data Structures and Algorithms I}
\newcommand{\ModuleCode}{CS201-S27}
\newcommand{\Programme}{Foundation programme}
\newcommand{\Level}{Semester 2}
\newcommand{\Semester}{Spring 2027}
\newcommand{\Credits}{TODO(owner)}
\newcommand{\ModuleLeader}{Dr. Ali Khudiyev}
\newcommand{\Email}{ali.khudiyev@wosm.academy}
\newcommand{\OfficeHours}{Monday--Friday, 10:00--17:00 (Baku time, UTC+4)}
\newcommand{\Prerequisites}{CS101-F26 Introduction to Informatics; M101-F26 Discrete Structures}

% ---------- Styling ------------------------------------------
\definecolor{accent}{RGB}{0,70,130}
\titleformat{\section}{\large\bfseries\color{accent}}{\thesection.}{0.5em}{}
\titlespacing*{\section}{0pt}{1.2em}{0.5em}
\setlist{itemsep=2pt, topsep=4pt}
\setlength{\parindent}{0pt}
\setlength{\parskip}{4pt}

\pagestyle{fancy}
\fancyhf{}
\lhead{\small \ModuleCode\ -- \ModuleTitle}
\rhead{\small Module Description}
\cfoot{\small \thepage}
\renewcommand{\headrulewidth}{0.4pt}

% =============================================================
\begin{document}

\begin{center}
  {\LARGE\bfseries\color{accent} \ModuleTitle}\\[4pt]
  {\large Module Description}
\end{center}
\vspace{6pt}

% ---------- Summary table (not shown on the website: the course page
% shows these details from the macros above) --------------------------
\begin{tabularx}{\textwidth}{>{\bfseries}l X >{\bfseries}l X}
\toprule
Module code   & \ModuleCode    & Credits      & \Credits \\
Programme     & \Programme     & Level        & \Level \\
Semester      & \Semester      & Prerequisites & \Prerequisites \\
Module leader & \ModuleLeader  & Email        & \href{mailto:\Email}{\Email} \\
Office hours  & \multicolumn{3}{l}{\OfficeHours} \\
\bottomrule
\end{tabularx}

% ---------- Sections (shown on the website) ------------------
\section{Module Overview}
Data Structures and Algorithms I is the first of two modules on the design and
analysis of algorithms. It teaches you to reason precisely about two questions that
every program raises: \emph{is it correct?} and \emph{how much time and memory does
it need?} You will study the fundamental data structures and the classic algorithms
for sorting and searching, and learn to choose between them by comparing the cost of
their operations.

The module builds directly on Semester 1. Programming is done in C. C exposes how
data structures are laid out in memory, which makes the cost of each operation
visible; the first units introduce C to students who learnt to program in Java in
Introduction to Informatics. Proofs by induction, asymptotic reasoning and basic
combinatorics come from Discrete Structures. Data Structures and Algorithms II
(Semester 3) continues with graph algorithms and general design techniques.

\section{Aims}
\begin{itemize}
  \item Introduce rigorous methods for proving algorithms correct and for analysing their running time and memory use.
  \item Give you a working toolkit of standard data structures and algorithms, and the judgement to choose between them.
  \item Develop the habit of thinking about data structures as abstract data types with separately analysed implementations.
  \item Introduce programming in C, and strengthen your programming by implementing, testing and measuring the data structures you study, with explicit control over memory.
\end{itemize}

\section{Learning Outcomes}
After completing this course, you will be able to:
\begin{enumerate}[label=LO\arabic*.]
  \item Analyse the running time and memory use of iterative and recursive algorithms using asymptotic notation, and solve recurrences by substitution, recursion trees and the master theorem.
  \item Prove the correctness of algorithms using loop invariants and induction.
  \item Implement and use the standard abstract data types (lists, stacks, queues, priority queues, sets and maps) in C with array-based and linked representations, managing memory explicitly, and analyse dynamic arrays with amortised analysis.
  \item Explain, implement and compare comparison-based and linear-time sorting algorithms, and explain the $\Omega(n \log n)$ lower bound for comparison sorting.
  \item Design, implement and analyse hash tables and balanced search trees, and choose an appropriate data structure for a given problem.
\end{enumerate}

\section{Syllabus}
\begin{tabularx}{\textwidth}{c X}
\toprule
\textbf{Unit} & \textbf{Topic} \\
\midrule
1  & Introduction: algorithms, problems and the RAM model \\
2  & C for Java programmers I: types, control flow, functions, pointers and arrays \\
3  & C for Java programmers II: strings, structs, dynamic memory and the toolchain \\
4  & Asymptotic notation and the growth of functions \\
5  & Correctness: loop invariants and proofs by induction \\
6  & Recursion and recurrences: substitution, recursion trees, master theorem \\
7  & Arrays, dynamic arrays and amortised analysis \\
8  & Linked lists, stacks and queues \\
9  & Abstract data types in C: header files, opaque types and function pointers \\
10 & Elementary sorting; merge sort \\
11 & Quicksort and randomised analysis \\
12 & Heaps, heapsort and priority queues \\
13 & Sorting lower bound; counting sort, radix sort and bucket sort \\
14 & Searching: binary search; selection and order statistics \\
15 & Hashing: hash functions, chaining and open addressing \\
16 & Binary search trees \\
17 & Balanced search trees: AVL and red-black trees \\
18 & Course Review \& Exam \\
\bottomrule
\end{tabularx}

\section{Teaching and Learning Methods}
The module is self-paced. All materials are provided on Moodle.
Each unit has a lecture, exercises on analysis and proofs, and a programming task in
C in which you implement and test the data structure or algorithm of the unit.

\begin{tabularx}{0.6\textwidth}{X r}
\toprule
\textbf{Activity} & \textbf{Hours} \\
\midrule
Lectures            & 34 \\
Labs / Tutorials    & 40 \\
Independent study   & 126 \\
\midrule
\textbf{Total}      & \textbf{200} \\
\bottomrule
\end{tabularx}

\section{Assessment}
Your knowledge will be assessed throughout the course with quizzes, programming
assignments and the final exam. The quizzes and the exam check whether you can
analyse algorithms, prove them correct and select suitable data structures. The
programming assignments check whether you can implement data structures and
algorithms correctly and efficiently.

\begin{tabularx}{\textwidth}{X c c c}
\toprule
\textbf{Component} & \textbf{Weight} & \textbf{Due} & \textbf{LOs} \\
\midrule
Quizzes & 20\% & Throughout & LO1--LO5 \\
Programming assignments & 40\% & Throughout & LO3--LO5 \\
Final exam & 40\% & End of module & LO1--LO5 \\
\midrule
\textbf{Total} & \textbf{100\%} & & \\
\bottomrule
\end{tabularx}

\section{Reading List}
\textbf{Essential}
\begin{itemize}
  \item T.\ H.\ Cormen, C.\ E.\ Leiserson, R.\ L.\ Rivest, C.\ Stein: \textit{Introduction to Algorithms}, 4th ed., MIT Press, 2022.
\end{itemize}
\textbf{Recommended}
\begin{itemize}
  \item R.\ Sedgewick: \textit{Algorithms in C, Parts 1--4: Fundamentals, Data Structures, Sorting, Searching}, 3rd ed., Addison-Wesley, 1998.
  \item B.\ W.\ Kernighan, D.\ M.\ Ritchie: \textit{The C Programming Language}, 2nd ed., Prentice Hall, 1988.
  \item K.\ N.\ King: \textit{C Programming: A Modern Approach}, 2nd ed., W.\ W.\ Norton, 2008.
  \item S.\ S.\ Skiena: \textit{The Algorithm Design Manual}, 3rd ed., Springer, 2020.
\end{itemize}

\section{Materials and Licence}
Course materials are published under the
\href{https://creativecommons.org/licenses/by-sa/4.0/}{Creative Commons Attribution-ShareAlike 4.0}
licence unless stated otherwise. See \href{https://wosm.academy/licence/}{wosm.academy/licence} for details.

\section{Policies}
\begin{itemize}
  \item \textbf{Attendance:} Attendance is not mandatory. The module is self-paced, so you can study the materials when it suits you.
  \item \textbf{Late submissions:} Work submitted after the deadline is penalised.
  \item \textbf{Academic integrity:} All submitted work must be your own.
        See the \href{https://wosm.academy/code-of-conduct/}{Code of Conduct}.
\end{itemize}

\end{document}
